%%=============================================================================
%% Fase 5: Testen
%%=============================================================================

\chapter{\IfLanguageName{dutch}{Testen}{Testing}}%
\label{ch:testen}

Dit hoofdstuk beschrijft de uitvoering van de testfase op het gemeentelijk
kerkhof van Sint-Gillis-Waas. De testopzet zoals gedefinieerd in de methodologie
werd gevolgd, met een technische evaluatie van alle kernfunctionaliteiten van de
applicatie en een observationele meting van de navigatietijd en gebruiksgemak.

%%─────────────────────────────────────────────────────────────────────────────
\section{Testopzet}
%%─────────────────────────────────────────────────────────────────────────────

De tests werden uitgevoerd op het gemeentelijk kerkhof van Sint-Gillis-Waas,
de testlocatie die werd geanalyseerd in hoofdstuk~\ref{ch:voorbereiding-analyse}.
De omstandigheden tijdens de test waren als volgt:

\begin{itemize}
    \item \textbf{Weersomstandigheden:} helder en zonnig
    \item \textbf{Tijdstip:} overdag
    \item \textbf{Toestel:} iPhone met iOS, ARKit-compatibel
    \item \textbf{Testpersonen:} % TODO: aantal testpersonen invullen na uitvoering
\end{itemize}

Zoals beschreven in de methodologie werden drie testscenario's gedefinieerd.
In deze testronde werd \textbf{Scenario C} (kaartweergave met AR) volledig
doorlopen. % TODO: scenario A en B toevoegen indien uitgevoerd met extra testpersonen.

%%─────────────────────────────────────────────────────────────────────────────
\section{Scenario C: Kaartweergave met AR}
%%─────────────────────────────────────────────────────────────────────────────

\subsection{Verloop van de test}

De testpersoon opende de applicatie aan de ingang van het kerkhof en zocht via
de zoekknop een specifiek graf op uit de mockdataset. Na selectie van het graf
startte de AR-navigatie onmiddellijk: de 3D-navigatiepijl verscheen in de
cameraweergave en de minikaart toonde de positie van het geselecteerde graf.
De testpersoon volgde de aanwijzingen van de applicatie tot het geselecteerde
graf werd bereikt.

De totale navigatietijd van de ingang tot aan het graf bedroeg
\textbf{approximately 5 minuten}.
% TODO: navigatietijden van extra testpersonen toevoegen

\subsection{AR-navigatiepijl}

De navigatiepijl functioneerde correct gedurende de volledige test. De pijl
wees op elk moment in de richting van het geselecteerde graf en roteerde vloeiend
mee wanneer de testpersoon van richting veranderde. Er werden geen fouten of
onverwachte sprong-bewegingen waargenomen.

Een relevante observatie betreft de aard van de navigatiebegeleiding. De pijl
wijst steeds in een rechte lijn naar het geselecteerde graf, ongeacht de
aanwezige paden of obstakels. In de praktijk betekent dit dat de gebruiker
zelf moet bepalen welk pad hij neemt om de bestemming te bereiken. Een ideale
navigatieoplossing zou de gebruiker langs de bestaande looppaden leiden, met
instructies zoals ``neem links, dan rechts''. Dit vereist echter GIS-paddata
van het kerkhof, zoals uitgebreid besproken in sectie~\ref{sec:routelijn}.
De rechte pijlnavigatie is dan ook geen technische tekortkoming van de
AR-component, maar een directe consequentie van de ontbrekende paddata.

\subsection{Grafmarker}

De world-locked grafmarker verscheen op het moment dat de testpersoon zich op
circa \textbf{20 meter} van de geselecteerde grafpositie bevond. Op die afstand
was de verticale lichtbalk duidelijk zichtbaar boven de omgeving, wat het
beoogde doel van de marker als visueel baken bevestigt.

Wat betreft de positioneringsnauwkeurigheid: de marker landde in de nabijheid
van de correcte grafsteen. Een exacte plaatsing is echter inherent beperkt
door de aard van de mockdata: de coördinaten werden handmatig ingevoerd op
basis van Google Maps en zijn benaderingen van de werkelijke grafposities.
Voor productiewaardige nauwkeurigheid zijn officieel ingemeten coördinaten
vereist, zoals die beschikbaar zijn in het WebBGP-systeem van de gemeente.

\subsection{ARKit tracking}

Het trackingstatusslabel toonde \textbf{onmiddellijk ``AR Active''} na het
opstarten van de applicatie. De initialisatietijd bedroeg ongeveer
\textbf{1 seconde}. Gedurende de volledige test bleef de trackingstatus stabiel,
zonder degradatie naar de statussen ``Insufficient Features'' of ``Slow Down''.
Dit is een positief resultaat voor een zonnige buitenomgeving.

\subsection{Fallback-banner}

De automatische fallback-banner werd tijdens de normale navigatie niet
geactiveerd, wat correct is: de AR-tracking was stabiel. De werking van de
banner werd apart getest door de cameralens gedurende enkele seconden af te
dekken. Na de geconfigureerde vertraging van 3 seconden verscheen de banner
correct met de opties om de kaart te openen of de banner te sluiten. Bij het
vrijgeven van de camera herstelde de tracking naar normaal en verdween de
banner automatisch.

\subsection{Zoekfunctie en minikaart}

De zoekfunctie vond de correcte grafrecords bij het intypen van een naam. De
minikaart toonde de gebruikerspositie en het geselecteerde graf correct. De
volledigschermanimatie van de minikaart werkte vloeiend. Het afstandslabel
toonde een nauwkeurige benadering van de afstand en werkte continu bij.

%%─────────────────────────────────────────────────────────────────────────────
\section{Samenvatting van de testresultaten}
%%─────────────────────────────────────────────────────────────────────────────

Tabel~\ref{tab:testresultaten} geeft een overzicht van de meetindicatoren en
de bijhorende resultaten.

\begin{table}[h]
    \centering
    \caption[Overzicht testresultaten]{Overzicht van de meetindicatoren en
        resultaten van de technische evaluatie (Scenario C).}
    \label{tab:testresultaten}
    \renewcommand{\arraystretch}{1.4}
    \begin{tabular}{p{4.5cm} p{7.5cm}}
        \hline
        \textbf{Indicator} & \textbf{Resultaat} \\
        \hline
        Navigatietijd & $\pm$ 5 minuten van ingang tot graf \\
        \hline
        Pijlrichting correct & Ja, continu correct gedurende de volledige test \\
        \hline
        Marker zichtbaar vanaf & $\pm$ 20 meter \\
        \hline
        GPS-positionering marker & Benaderd correct; nauwkeurigheid beperkt door mockdata \\
        \hline
        ARKit tracking & Stabiel; ``AR Active'' na $\pm$ 1 seconde \\
        \hline
        Fallback-banner & Correct geactiveerd bij gesimuleerde tracking-degradatie \\
        \hline
        Zoekfunctie & Correct; graven gevonden op naam \\
        \hline
        Minikaart positie & Correct; gebruikerspositie en grafannotatie zichtbaar \\
        \hline
        Afstandslabel & Nauwkeurig en continu bijgewerkt \\
        \hline
        Weersomstandigheden & Zonnig; geen tracking-problemen vastgesteld \\
        \hline
    \end{tabular}
\end{table}

% TODO: tabel uitbreiden met resultaten van extra testpersonen